\section{Переобкачка}

Переобкачка вынесена в отдельную программу, \texttt{src/recrawler.py}. Она ничего не
скачивает \emph{в базу}: её задача --- определить, изменился ли документ, и если
да, поставить его URL в очередь. Дальше документ заберёт обычный робот.

Программа берёт из базы документы, которые давно не проверялись (порог задаётся
в секции \texttt{recrawl} конфига), и по каждому идёт в два шага --- от дешёвого
к дорогому.

Периодичность обеспечивает системный планировщик, а не сама программа: например,
строка в \texttt{crontab}
\begin{verbatim}
0 3 * * *  cd ~/MAI-IR/labs/lab-02 && make recrawl
\end{verbatim}
запускает проверку каждую ночь. Порог \texttt{min\_age\_hours} не даёт проверять
один документ чаще раза в сутки, сколько бы раз программу ни запустили, а
\texttt{batch} ограничивает число запросов за один запуск.

\textbf{Шаг 1. Условный запрос.} Если у документа сохранены \texttt{ETag} или
\texttt{Last-Modified}, запрос отправляется с заголовками
\texttt{If-None-Match} и \texttt{If-Modified-Since}. Ответ \texttt{304 Not
Modified} означает, что документ точно не менялся, --- причём тело страницы
сервер при этом не передаёт вовсе.

\textbf{Шаг 2. Сравнение отпечатка.} Если сервер условные запросы не
поддерживает, страница скачивается, из неё выделяется текст и считается
отпечаток SHA-256. Он сравнивается с сохранённым в базе.

Сравнивается именно \emph{выделенный текст}, а не \enquote{сырой} HTML. Это
принципиально: в HTML на каждой перезагрузке меняются счётчики, идентификаторы
сессий, рекламные блоки и метки времени, и по нему почти любой документ выглядел
бы изменившимся. Отпечаток текста меняется только тогда, когда изменилось
содержание.

Изменившийся URL попадает в очередь со статусом \texttt{pending} и приоритетом
1. Приоритет важен: очередь отдаёт такие адреса раньше новых, поэтому обновление
уже известных документов не ждёт, пока разгребётся весь фронт обкачки. При
повторной обкачке запись в \texttt{documents} обновляется, а счётчик
\texttt{revision} увеличивается на единицу.

\subsection*{Проверка работы}

Цикл проверялся на настоящих данных дважды.

\textbf{Первая проверка, 7 октября.} Через неделю после демонстрационного
прогона (раздел \enquote{Результаты}) все 100 документов тогдашней базы стали
старше суток, и переобкачка проверила их целиком:

\begin{verbatim}
[изменился] https://ru.wikipedia.org/wiki/7_ящиков
[изменился] https://ru.wikipedia.org/wiki/9_песен
...
Время работы:              113.7 с
Проверено документов:      100
  ответ 304 Not Modified:  68
  текст совпал:            24
  изменились:              8
В очередь на переобкачку добавлено: 8
\end{verbatim}

Две трети документов отсеялись первым, дешёвым шагом: Википедия поддерживает
условные запросы, и на \texttt{304} тело страницы не передаётся. У 24 документов
сервер отдал страницу целиком, но выделенный текст не изменился, --- это ровно
тот случай, ради которого сравнивается текст, а не HTML: по HTML все 24
считались бы изменившимися.

Восемь изменений проверены вручную, и все они настоящие правки статей. В
\enquote{7 ящиках} появились разделы о критике и наградах. Из \enquote{9 песен}
и \enquote{8 женщин} редакторы удалили разделы с плашкой \enquote{В разделе не
хватает ссылок на источники}. В остальных поправлены отдельные слова, цифры и
типографика.

Запущенный следом робот взял все восемь адресов раньше новых и поднял им
\texttt{revision} с 1 до 2:

\begin{verbatim}
[updated ]    2214 симв.  https://ru.wikipedia.org/wiki/7_ящиков
[updated ]    2766 симв.  https://ru.wikipedia.org/wiki/9_песен
...
\end{verbatim}

Эта проверка нашла и ошибку. Сначала робот перекачал один адрес и переключился
на новые страницы Кино-Театр.Ру: после запроса к Википедии её хост был занят
паузой вежливости, а выбор по готовности хоста ставил готовый источник выше
приоритета переобкачки. Теперь очередь сначала ищет адреса, ждущие переобкачки,
во всех источниках, и только потом --- новые; на это добавлен тест.

\textbf{Вторая проверка, 19 октября,} --- уже на полном корпусе. Переобкачка взяла
100 документов, которые дольше всех не проверялись (это те же статьи, что и в
первый раз):

\begin{verbatim}
Проверено документов:      100
  ответ 304 Not Modified:  0
  текст совпал:            97
  изменились:              3
\end{verbatim}

На этот раз ни один документ не отсеялся условным запросом. Причина не в
роботе: 14--15 октября Википедия обновила у всех этих статей дату
\texttt{Last-Modified}, хотя текст не менялся, --- по-видимому, массово правились
шаблоны или категории. Сервер честно отдал страницы целиком, и второй шаг,
сравнение отпечатков текста, отсеял 97 неизменившихся. Это хорошо показывает,
зачем нужны оба шага: условный запрос дешёвый, но полагаться на одну дату
изменения нельзя. Новые даты сохранены, так что при следующей проверке эти
документы ответят \texttt{304}. Три изменившиеся статьи робот перекачал первыми.

\pagebreak
